% Composition study only — force mechanism is deliberately the visual pivot.
\documentclass[border=4pt]{standalone}
\usepackage{polymer-paper-preamble}
\usetikzlibrary{arrows.meta,calc,positioning}

\begin{document}
\resizebox{178mm}{!}{%
\begin{tikzpicture}[
  every node/.style={font=\sffamily\fontsize{5.4}{6.5}\selectfont},
  panelLetter/.style={font=\sffamily\bfseries\fontsize{8}{9.6}\selectfont,text=cGray!92!black},
  panelTitle/.style={font=\sffamily\bfseries\fontsize{6.5}{7.8}\selectfont,text=cGray!86!black},
  label/.style={font=\sffamily\fontsize{5.2}{6.2}\selectfont,text=cGray!85!black},
  mute/.style={font=\sffamily\itshape\fontsize{5.05}{6.0}\selectfont,text=cGray!62!black},
  apparatus/.style={draw=cGray!76!black,fill=cGray!18,line width=.75pt},
  wire/.style={draw=cGray!78!black,line width=.75pt,line cap=round},
  beamOuter/.style={cAmber!74!black,line width=2.35mm,line cap=round},
  beamInner/.style={cAmber!10,line width=1.52mm,line cap=round},
  leader/.style={cGray!58!black,line width=.40pt,line cap=round},
  stageArrow/.style={-{Stealth[length=3.8pt,width=2.7pt]},cGray!64!black,line width=.65pt},
  force/.style={-{Stealth[length=4.7pt,width=3.4pt]},cRed!86!black,line width=1.00pt},
  baseline/.style={-{Stealth[length=3.5pt,width=2.5pt]},cGray!62!black,line width=.60pt},
  q/.style={circle,fill=cRed!80!black,draw=cRed!95!black,line width=.26pt,minimum size=1.22mm,inner sep=0pt},
  terminal/.style={circle,draw=cGray!76!black,fill=white,line width=.54pt,minimum size=2.8pt,inner sep=0pt},
  axis/.style={-{Stealth[length=3.2pt,width=2.3pt]},cGray!72!black,line width=.60pt},
]

% Preparation and isolation are compact process anchors, not co-equal visual subjects.
\begin{scope}[shift={(0,5.15)}]
  \node[panelLetter,anchor=west] at (.10,2.27) {a};
  \node[panelTitle,anchor=west] at (.45,2.27) {actuation charge};
  \node[mute,anchor=west] at (.45,1.94) {field-on actuation};
  \fill[apparatus] (1.25,1.38) rectangle (2.25,1.65);
  \draw[wire] (1.75,1.38)--(1.75,1.93);
  \node[mute,anchor=west] at (2.00,1.71) {clip: GND};
  \draw[beamOuter] (1.75,1.38) .. controls (1.75,.95) and (2.12,.45) .. (2.78,.28);
  \draw[beamInner] (1.75,1.38) .. controls (1.75,.95) and (2.12,.45) .. (2.78,.28);
  \fill[apparatus] (4.10,.18) rectangle (4.38,1.65);
  \draw[leader] (4.24,1.65)--(4.24,1.92);
  \node[label,text=cRed!82!black,anchor=south] at (4.24,1.92) {$+5\,\mathrm{kV}$};
  \node[q] at (1.83,.98) {}; \node[q] at (2.18,.53) {};
  \draw[force] (1.89,.92)--(3.45,.92);
  \node[label,text=cRed!84!black,anchor=south] at (2.68,.98) {attraction};
\end{scope}

\begin{scope}[shift={(5.20,5.15)}]
  \node[panelLetter,anchor=west] at (.10,2.27) {b};
  \node[panelTitle,anchor=west] at (.45,2.27) {source-off isolation};
  \node[mute,anchor=west] at (.45,1.94) {lead lifted; clip floating};
  \fill[apparatus] (1.25,1.38) rectangle (2.25,1.65);
  \draw[wire] (1.75,1.38)--(1.75,1.84);
  \node[mute,anchor=west] at (2.00,1.69) {clip mounted};
  \draw[beamOuter] (1.75,1.38) .. controls (1.75,.94) and (1.90,.48) .. (2.36,.24);
  \draw[beamInner] (1.75,1.38) .. controls (1.75,.94) and (1.90,.48) .. (2.36,.24);
  \fill[apparatus] (4.10,.18) rectangle (4.38,1.65);
  \node[q] at (1.81,.95) {}; \node[q] at (2.05,.48) {};
  \draw[wire] (.75,.45)--(1.18,.45); \node[terminal] at (1.18,.45) {};
  \draw[wire] (1.76,.70)--(2.19,.70); \node[terminal] at (1.76,.70) {};
  \draw[baseline] (1.47,.50)--(1.47,.65);
  \node[mute,anchor=north] at (1.47,.32) {source OFF};
\end{scope}
\draw[stageArrow] (4.66,6.22)--(5.03,6.22);
\draw[stageArrow] (7.75,5.03)--(7.75,4.58);
\node[mute,anchor=west] at (7.88,4.80) {reverse drive};

% C occupies the main mechanism field: the conditional result is read before the plot.
\begin{scope}[shift={(0,.00)}]
  \node[panelLetter,anchor=west] at (.10,4.40) {c};
  \node[panelTitle,anchor=west] at (.45,4.40) {force balance after isolation};
  \node[mute,anchor=west] at (.45,4.05) {same mounted film; conditional reverse bend};
  \fill[apparatus] (2.05,3.28) rectangle (3.38,3.62);
  \coordinate (c-root) at (2.72,3.28);
  \draw[wire] (c-root)--(2.72,3.73);
  \node[mute,anchor=west] at (3.12,3.76) {electrically floating clip};
  \draw[beamOuter] (c-root) .. controls (2.72,2.54) and (2.08,1.23) .. (1.02,.62);
  \draw[beamInner] (c-root) .. controls (2.72,2.54) and (2.08,1.23) .. (1.02,.62);
  \fill[apparatus] (7.62,.48) rectangle (7.98,3.62);
  \draw[leader] (7.80,3.62)--(7.80,3.90);
  \node[label,text=cRed!82!black,anchor=south] at (7.80,3.90) {$-5\,\mathrm{kV}$};
  \node[q] at (2.55,2.50) {}; \node[q] at (1.93,1.47) {}; \node[q] at (1.10,.69) {};
  \node[label,text=cRed!80!black,anchor=east] at (1.48,2.82) {trapped $q_{\mathrm{tr}}$};
  \draw[leader,cRed!58!black] (1.54,2.76)--(2.55,2.50);
  \draw[baseline] (2.72,2.66)--(4.78,2.66);
  \node[label,anchor=south west,text=cGray!72!black] at (4.86,2.70) {Maxwell attraction};
  \draw[force] (1.93,1.47)--(.30,1.47);
  \node[label,text=cRed!84!black,anchor=south east] at (1.50,1.53) {Coulomb};
  \node[label,anchor=north] at (4.95,.52) {reverse bend if $|F_{\mathrm C}|>|F_{\mathrm M}|$};
\end{scope}

% D is a compact outcome panel, connected by the condition rather than a duplicated device cartoon.
\begin{scope}[shift={(9.30,.00)}]
  \node[panelLetter,anchor=west] at (.10,4.40) {d};
  \node[panelTitle,anchor=west] at (.45,4.40) {continuous response};
  \node[mute,anchor=west] at (.45,4.05) {qualitative bend trace};
  \draw[axis] (.72,.72)--(.72,3.25);
  \draw[axis] (.72,1.90)--(7.85,1.90);
  \node[mute,rotate=90,anchor=south] at (.30,1.96) {bend angle};
  \node[mute,anchor=north east] at (7.85,.60) {time};
  \node[mute,anchor=east] at (.64,2.83) {$+\theta$};
  \node[mute,anchor=east] at (.64,1.08) {$-\theta$};
  \draw[cBlue!82!black,line width=1.04pt]
    (.86,1.90) .. controls (.98,2.04) and (1.25,2.68) .. (1.55,2.80)
    -- (3.23,2.80)
    .. controls (3.42,2.80) and (3.52,2.66) .. (3.63,1.90)
    .. controls (3.77,1.17) and (3.98,.94) .. (4.46,.97)
    .. controls (5.14,1.03) and (5.77,1.44) .. (6.42,1.66)
    .. controls (6.92,1.82) and (7.38,1.87) .. (7.65,1.88);
  \draw[cBlue!55!white,dashed,line width=.80pt] (7.65,1.88)--(7.85,1.90);
  \draw[cGray!45!black,dashed,line width=.50pt] (3.23,.93)--(3.23,2.80);
  \node[mute,anchor=south] at (3.23,3.32) {source OFF; floating};
  \node[label,anchor=south] at (2.36,2.89) {positive plateau};
  \node[label,text=cRed!84!black,anchor=south west] at (3.70,2.22) {reversed drive};
  \node[mute,anchor=south west] at (5.78,1.66) {slow recovery};
  \draw[cGray!70!black,line width=.45pt] (.86,1.90)--(.86,1.70);
  \node[mute,anchor=north] at (.86,1.63) {$t=0$};
\end{scope}
\end{tikzpicture}%
}
\end{document}
